\section{答题模板}

本章按题型和知识模块整理答题规范模板，帮助考生建立稳定的解题框架。

\subsection{选择题模板}

\subsubsection{运动学图象分析模板}

\textbf{问题特征}：给出 $v-t$、$x-t$、$a-t$ 图象，判断运动性质、位移、加速度等。

\textbf{解题步骤}：
\begin{enumerate}
  \item 明确坐标轴物理意义（横轴、纵轴代表的物理量）
  \item 读图象要素：
    \begin{itemize}
      \item 斜率：$v-t$ 图斜率 = 加速度，$x-t$ 图斜率 = 速度
      \item 面积：$v-t$ 图与横轴围成的面积 = 位移
      \item 截距：初始值
    \end{itemize}
  \item 分段分析，注意图象转折点对应的运动状态变化
  \item 结合牛顿定律或运动学公式进一步判断
\end{enumerate}

\textbf{常见陷阱}：
\begin{itemize}
  \item $x-t$ 图不是运动轨迹
  \item $v-t$ 图中面积在横轴下方表示负方向位移
  \item 图象交点含义：$v-t$ 图中交点表示速度相等而非相遇
\end{itemize}

\subsubsection{受力分析模板}

\textbf{问题特征}：物体在多个力作用下处于平衡或加速状态。

\textbf{解题步骤}：
\begin{enumerate}
  \item 明确研究对象（整体法或隔离法）
  \item 画受力分析图：
    \begin{itemize}
      \item 先画重力（方向竖直向下）
      \item 再画弹力（垂直于接触面）
      \item 最后画摩擦力（沿接触面，与相对运动/趋势方向相反）
    \end{itemize}
  \item 建立坐标系（通常沿加速度方向建轴）
  \item 列平衡方程（$\sum F_x=0,\sum F_y=0$）或牛顿第二定律（$\sum F = ma$）
\end{enumerate}

\textbf{常用技巧}：
\begin{itemize}
  \item 整体法：系统内各物体相对静止时，可视为整体
  \item 隔离法：求内力时必须隔离
  \item 正交分解法：尽量让多的力落在坐标轴上
\end{itemize}

\subsubsection{功能关系模板}

\textbf{问题特征}：涉及动能、势能、机械能变化，判断做功情况。

\textbf{解题步骤}：
\begin{enumerate}
  \item 明确研究对象和过程
  \item 分析各力做功：
    \begin{itemize}
      \item 重力、弹力做功 $\to$ 对应势能变化
      \item 摩擦力做功 $\to$ 产生内能
      \item 安培力做功 $\to$ 电能转化
    \end{itemize}
  \item 根据问题选择公式：
    \begin{itemize}
      \item 求速度/位移 $\to$ 动能定理 $W_{\text{合}} = \frac12 mv_2^2 - \frac12 mv_1^2$
      \item 判断机械能是否守恒 $\to$ 除重力（弹力）外无其他力做功
      \item 涉及摩擦生热 $\to$ $Q = f\cdot s_{\text{相}}$
    \end{itemize}
  \item 列方程求解
\end{enumerate}

\subsubsection{电路动态分析模板}

\textbf{问题特征}：滑动变阻器滑片移动或开关通断，判断各电表示数变化。

\textbf{常用方法——串反并同法}：
\begin{itemize}
  \item 变量增大时：
    \begin{itemize}
      \item 与变量串联的元件，其物理量变化趋势相反（串反）
      \item 与变量并联的元件，其物理量变化趋势相同（并同）
    \end{itemize}
  \item 例：滑片右移 $\to$ 电阻增大 $\to$ 总电流减小 $\to$ 路端电压增大
\end{itemize}

\textbf{常规分析路径}：
\[
\begin{array}{c}
\text{局部电阻变化} \to \text{总电阻变化} \to \text{总电流变化} \\
\to \text{路端电压变化} \to \text{支路电流和电压变化}
\end{array}
\]

\subsection{实验题模板}

\subsubsection{力学实验——纸带处理模板}

\textbf{标准步骤}：
\begin{enumerate}
  \item 确定计数点间的位移 $x_1,x_2,\cdots,x_n$
  \item 求某点瞬时速度：$v_n = \dfrac{x_n + x_{n+1}}{2T}$（$T$ 为计数点时间间隔）
  \item 求加速度：$a = \dfrac{(x_4+x_5+x_6) - (x_1+x_2+x_3)}{9T^2}$（逐差法）
  \item 分析误差：存在阻力 $\to$ 加速度偏小
\end{enumerate}

\subsubsection{电学实验——伏安法测电阻模板}

\textbf{器材选择}：
\begin{enumerate}
  \item 电源：电动势略大于被测电路所需电压
  \item 电压表：量程略大于被测电压
  \item 电流表：量程略大于被测电流
  \item 滑动变阻器：阻值与待测电阻可比，便于调节
\end{enumerate}

\textbf{电路选择}：
\begin{itemize}
  \item 内接法：$\dfrac{R_x}{R_A} > \dfrac{R_V}{R_x}$ 时选用（大电阻）
    \begin{itemize}
      \item 测量值 $= R_x + R_A$，偏大
    \end{itemize}
  \item 外接法：$\dfrac{R_x}{R_A} < \dfrac{R_V}{R_x}$ 时选用（小电阻）
    \begin{itemize}
      \item 测量值 $= \dfrac{R_xR_V}{R_x+R_V}$，偏小
    \end{itemize}
\end{itemize}

\textbf{滑动变阻器接法}：
\begin{itemize}
  \item 限流式：节能，适用于滑动变阻器阻值与负载可比时
  \item 分压式：电压从零开始可调，适用于滑动变阻器阻值较小时
  \item \textbf{必选分压式的三种情况}：电压从零开始、滑动变阻器阻值远小于负载、要求多测几组数据
\end{itemize}

\subsubsection{电学实验——测电源电动势和内阻模板}

\textbf{原理}：$U = E - Ir$

\textbf{电路连接}：电流表测干路电流，电压表测路端电压

\textbf{数据处理}：
\begin{enumerate}
  \item 改变滑动变阻器阻值，测多组 $U$、$I$
  \item 作 $U-I$ 图象
  \item 纵截距 $= E$，斜率的绝对值 $= r$
\end{enumerate}

\textbf{误差分析}：
\begin{itemize}
  \item 电流表外接法：$E_{\text{测}} < E_{\text{真}}$，$r_{\text{测}} < r_{\text{真}}$
  \item 电流表内接法：$E_{\text{测}} = E_{\text{真}}$，$r_{\text{测}} > r_{\text{真}}$
\end{itemize}

\subsection{计算题模板}

\subsubsection{力学综合——多过程问题模板}

\textbf{问题特征}：物体经历多个运动阶段（直线 $\to$ 曲线 $\to$ 碰撞等）。

\textbf{解题步骤}：
\begin{enumerate}
  \item \textbf{过程拆分}：按运动性质将全过程分为若干子过程
  \item \textbf{状态分析}：确定每个子过程的初末状态（速度、位置、能量）
  \item \textbf{选择规律}：
    \begin{itemize}
      \item 匀变速直线：运动学公式
      \item 圆周运动：牛顿第二定律 $F_n = m\frac{v^2}{r}$
      \item 平抛运动：分解为水平和竖直方向
      \item 碰撞/爆炸：动量守恒（系统合外力为零时）
    \end{itemize}
  \item \textbf{衔接条件}：前一过程的末状态 = 后一过程的初状态
  \item \textbf{联立求解}
\end{enumerate}

\textbf{模板示例}：
\[
\begin{aligned}
&\text{过程 1（匀加速）：} \quad v_1^2 = 2a_1x_1 \\
&\text{过程 2（圆周）：} \quad mg + N = m\frac{v_2^2}{R} \\
&\text{过程 3（平抛）：} \quad x = v_2t,\; h = \frac12 gt^2
\end{aligned}
\]

\subsubsection{力学综合——板块模型模板}

\textbf{问题特征}：滑块在木板上滑动，涉及相对运动、摩擦生热。

\textbf{解题步骤}：
\begin{enumerate}
  \item 分别对滑块和木板进行受力分析
  \item 由牛顿第二定律求各自的加速度
  \item 分析共速条件：
    \begin{itemize}
      \item 设 $t$ 时刻共速：$v_{\text{滑}} = v_{\text{板}}$
      \item 代入 $v = v_0 + at$ 求 $t$
    \end{itemize}
  \item 求相对位移：$s_{\text{相}} = s_{\text{滑}} - s_{\text{板}}$
  \item 摩擦生热：$Q = \mu mg\cdot s_{\text{相}}$
  \item 判断滑块是否滑离木板
\end{enumerate}

\subsubsection{力学综合——弹簧模型模板}

\textbf{问题特征}：涉及弹簧的压缩/伸长、弹性势能变化。

\textbf{解题步骤}：
\begin{enumerate}
  \item 分析弹簧状态（原长、压缩、伸长）和形变量
  \item 弹力 $F = kx$（注意方向）
  \item 弹性势能 $E_p = \frac12 kx^2$
  \item 机械能守恒（只有重力和弹力做功时）：
    \[
    \frac12 mv_1^2 + mgh_1 + \frac12 kx_1^2 = \frac12 mv_2^2 + mgh_2 + \frac12 kx_2^2
    \]
  \item 临界条件：速度最大时 $a=0$（合力为零），即 $kx = mg$
\end{enumerate}

\subsubsection{电学综合——带电粒子在电场中运动模板}

\textbf{问题特征}：带电粒子在匀强电场中加速或偏转。

\textbf{加速电场}：
\[
qU = \frac12 mv^2 \quad\Rightarrow\quad v = \sqrt{\frac{2qU}{m}}
\]

\textbf{偏转电场（类平抛）}：
\[
\begin{aligned}
a &= \frac{qE}{m} = \frac{qU}{md} \\
t &= \frac{l}{v_0} \\
y &= \frac12 at^2 = \frac{qUl^2}{2mdv_0^2} \\
\tan\theta &= \frac{v_y}{v_x} = \frac{qUl}{mdv_0^2}
\end{aligned}
\]

\textbf{关键结论}：
\begin{itemize}
  \item 粒子穿出电场时速度反向延长线交于极板中点
  \item 偏移量 $y$ 与 $q/m$ 成正比，与 $v_0^2$ 成反比
\end{itemize}

\subsubsection{电学综合——带电粒子在磁场中运动模板}

\textbf{问题特征}：带电粒子在匀强磁场中做匀速圆周运动。

\textbf{标准步骤——三步法}：

\textbf{第 1 步：确定圆心}
\begin{itemize}
  \item 方法一：两条速度垂线的交点
  \item 方法二：速度垂线与弦的中垂线交点
  \item 方法三：速度垂线与入射/出射点连线的中垂线交点
\end{itemize}

\textbf{第 2 步：求半径}
\[
\text{几何法：} R = \frac{L}{2\sin\theta},\quad
\text{公式法：} R = \frac{mv}{qB}
\]

\textbf{第 3 步：求时间}
\[
t = \frac{\theta}{2\pi}T = \frac{\theta}{2\pi}\cdot\frac{2\pi m}{qB} = \frac{m\theta}{qB}
\]
（$\theta$ 为轨迹圆心角，单位为弧度）

\textbf{常用几何关系}：
\begin{itemize}
  \item 入射速度与边界夹角 $=$ 出射速度与边界夹角（对称性）
  \item 粒子沿径向射入圆形磁场，必沿径向射出
  \item 轨迹半径 $R$、磁场圆半径 $r$、偏转角 $\theta$ 的几何关系
\end{itemize}

\subsubsection{电学综合——带电粒子在组合场中运动模板}

\textbf{问题特征}：粒子先后经过电场和磁场区域。

\textbf{常见模型}：
\begin{enumerate}
  \item \textbf{先电场后磁场}：
    \begin{itemize}
      \item 电场中加速/偏转 $\to$ 求进入磁场的速度
      \item 磁场中圆周运动 $\to$ 求半径和偏转角
      \item 出磁场后可能再做匀速直线运动
    \end{itemize}
  \item \textbf{先磁场后电场}：
    \begin{itemize}
      \item 磁场中圆周运动 $\to$ 求离开磁场的位置和速度方向
      \item 电场中类平抛 $\to$ 求偏转量
    \end{itemize}
  \item \textbf{交替场（回旋加速器）}：电场加速 $\to$ 磁场偏转 $\to$ 反复
\end{enumerate}

\subsubsection{电磁感应综合模板}

\textbf{问题特征}：导体棒在磁场中切割磁感线，涉及电路、安培力、能量转化。

\textbf{解题步骤}：
\begin{enumerate}
  \item \textbf{电路分析}：
    \begin{itemize}
      \item 切割导体等效为电源（内阻为导体棒电阻）
      \item 画出等效电路图
      \item 由 $E = Blv$ 或 $E = n\frac{\Delta\Phi}{\Delta t}$ 求感应电动势
    \end{itemize}
  \item \textbf{受力分析}：
    \begin{itemize}
      \item 安培力 $F_A = BIl = \frac{B^2l^2v}{R_{\text{总}}}$（阻碍相对运动）
      \item 列出牛顿第二定律方程
    \end{itemize}
  \item \textbf{运动分析}：
    \begin{itemize}
      \item 导体棒做加速度减小的加速运动（或减速）
      \item 最终匀速时安培力与外力平衡
      \item 最大速度 $v_m = \frac{F_{\text{外}}R_{\text{总}}}{B^2l^2}$
    \end{itemize}
  \item \textbf{能量分析}：
    \[
    W_{\text{外}} = Q_{\text{焦耳}} + \Delta E_k + \Delta E_p
    \]
    \[
    Q_{\text{焦耳}} = I^2Rt = \frac{E^2}{R_{\text{总}}}t
    \]
\end{enumerate}

\subsubsection{热学（选修 3--3）模板}

\textbf{问题特征}：气缸/活塞类问题，求气体压强、体积、温度变化。

\textbf{压强分析}：
\begin{itemize}
  \item 活塞平衡：$pS = p_0S + mg$（活塞水平时；竖直时考虑重力方向）
  \item 液柱平衡：$p = p_0 + \rho gh$（$h$ 为液面高度差）
\end{itemize}

\textbf{状态方程应用}：
\[
\frac{p_1V_1}{T_1} = \frac{p_2V_2}{T_2}
\]

\textbf{三步法}：
\begin{enumerate}
  \item 确定研究对象（一定质量的理想气体）
  \item 列出初末状态参量（$p_1,V_1,T_1$；$p_2,V_2,T_2$）
  \item 根据过程特点选择相应定律或状态方程
\end{enumerate}

\textbf{气体内能变化}：
\begin{itemize}
  \item 等温：$\Delta U = 0$，$Q = W$
  \item 等容：$W = 0$，$\Delta U = Q$
  \item 绝热：$Q = 0$，$\Delta U = W$
\end{itemize}

\subsubsection{光学（选修 3--4）模板}

\textbf{折射问题}：
\[
n = \frac{\sin i}{\sin r} = \frac{c}{v}
\]

\textbf{全反射问题}：
\[
\sin C = \frac{1}{n}
\]

\textbf{解题步骤}：
\begin{enumerate}
  \item 画光路图，标出入射角、折射角
  \item 判断是否发生全反射：光从光密介质射向光疏介质，且 $i > C$
  \item 应用折射定律或全反射条件列方程
  \item 结合几何关系求解（三角形内角和、平行线性质等）
\end{enumerate}

\subsubsection{原子物理模板}

\textbf{能级跃迁}：
\[
h\nu = E_m - E_n \quad (m > n)
\]
\begin{itemize}
  \item 大量氢原子处于 $n$ 能级时，可辐射 $C_n^2$ 种频率的光
  \item 从 $n$ 能级跃迁到基态，波长最长为 $n\to n-1$，最短为 $n\to1$
\end{itemize}

\textbf{核能计算}：
\[
\Delta E = \Delta m \times 931.5\,\text{MeV}
\]
（$\Delta m$ 以原子质量单位 u 为单位，$1\,\text{u} = 1.66\times10^{-27}\,\text{kg}$）
